% source: An algebraic approach to the Riemann-Roch theorem and the arithmetic theory of function fields (Athanasios Papaioannou), http://www.fen.bilkent.edu.tr/~franz/mat/Riemann-Roch.pdf
% pdf_page: 0013
% retrieved: 2026-07-26
% transcribed_by: page-transcriber (AI vision, no OCR)

\documentclass[11pt]{article}
\usepackage{amsmath,amssymb,amsthm}

% Small-caps theorem style matching the source.
\newtheoremstyle{scplain}{\topsep}{\topsep}{\itshape}{0pt}{\scshape}{.}{ }{\thmname{#1} \thmnote{#3}}
\theoremstyle{scplain}
\newtheorem*{theorem}{Theorem}
\newtheorem*{corollary}{Corollary}
\newtheorem*{conjecture}{Conjecture}

% Small-caps "Proof." to match the source.
\makeatletter
\renewenvironment{proof}[1][\proofname]{%
  \par\pushQED{\qed}\normalfont\topsep6\p@\@plus6\p@\relax
  \trivlist\item[\hskip\labelsep\scshape #1\@addpunct{.}]\ignorespaces}%
  {\popQED\endtrivlist\@endpefalse}
\makeatother
\renewcommand{\qedsymbol}{$\square$}

\DeclareMathOperator{\Div}{Div}

\begin{document}
\setlength{\emergencystretch}{2em}

\begin{flushright}$\square$\end{flushright}

As a final step before essentially completing the proof of Riemann-Roch, we
must also show the following

\begin{theorem}[1.26]
Let $D \in \Div(\mathbb{F})$ be an arbitrary divisor and $W=(\omega)$ a
canonical divisor of $\mathbb{F}/\mathbb{K}$. Then the mapping
$\mu : L(W-D) \longrightarrow \Omega_{\mathbb{F}}(D)$ defined by
$x \mapsto x\omega$ is an isomorphism of $\mathbb{K}$-vector spaces. In
particular,
\[
  i(D)=\dim(W-D).
\]
\end{theorem}

\begin{proof}
For $x \in L(W-D)$, we have
$(x\omega)=(x)+(\omega) \geq -(W-D)+W=D$, hence
$x\omega \in \Omega_{\mathbb{F}}(D)$. Therefore, $\mu$ maps $L(W-D)$ into
$\Omega_{\mathbb{F}}(D)$. Clearly, $\mu$ is linear and injective; in order to
show that $\mu$ is also surjective consider a Weil differential
$\omega_1 \in \omega_{\mathbb{F}}(D)$. By the fact that
$\Omega_{\mathbb{F}}$ has dimension $1$ over $\mathbb{F}$, there is
$x \in \mathbb{F}$ such that $\omega_1=x\omega$, and since
$(x)+W=(x\omega)=(\omega_1) \geq D$, we obtain $(x) \geq -(W-D)$, hence
$x \in L(W-D)$ and $\omega_1=\mu(x)$, as desired. Now, we have shown that
$\dim\Omega_{\mathbb{F}}(D)=i(D)$; this implies that $i(D)=\dim(W-D)$.
\end{proof}

Now we are ready to reach our goal for this section. Without any effort at all,
we can reach for the golden apple, namely

\begin{theorem}[1.27 (Riemann-Roch)]
Let $W$ be a canonical divisor of the function field
$\mathbb{F}/\mathbb{K}$. Then, for any divisor $D \in \Div(\mathbb{F})$,
\[
  \dim D=\deg D+1-g+\dim(W-D).
\]
\end{theorem}

\begin{proof}
This is an immediate consequence of the definition of $i(D)$ and the previous
theorem.
\end{proof}

A few obvious corollaries follow immediately:

\begin{corollary}[1.28]
Any canonical divisor $W$ has degree $2g-2$ and dimension $g$.
\end{corollary}

\begin{proof}
For $D=0$, the Riemann-Roch theorem yields $\dim W=g$ and then $D=W$ yields
$\deg W=2g-2$.
\end{proof}

\begin{corollary}[1.29]
If $D \in \Div(\mathbb{F})$ is such that $\deg D \geq 2g-1$, then
$\dim D=\deg D+1-g$.
\end{corollary}

\begin{proof}
We have $\dim D=\deg D+1-g+\dim(W-D)$ for any canonical divisor $W$; since
$\deg D \geq 2g-1$ and $\deg W=2g-2$ (by the previous corollary), we have
$\deg(W-D)<0$. This can only occur when $\dim(W-D)=0$, and this yields the
desired result.
\end{proof}

\section*{2\quad An arithmetic theory for function fields.}

Having developed some basic material about function fields, we are now going to
discuss briefly and informally the utility of this algebraic point of view on
function fields. The material we presented in the first section is usually
introduced in a geometric manner, function fields corresponding to algebraic
curves and their extensions to ramified covers of curves; however, we favoured
the algebraic point of view because it leads very naturally to an arithmetic
theory of function fields, which bears astonishing similarity to the classical
theory number fields.

The proofs, however, are significantly easier in the function field setting
than in the classical one. In fact, there are sweeping results that can be
proved fairly easily in the case of function fields but are still intractable
for number fields.

An example of such a theorem is the ABC conjecture of Oesterl\'e and Masser. In
the case of rational numbers $\mathbb{Q}$, the usual form of the conjecture has
as follows:

\begin{conjecture}[2.1 (ABC I, by Oesterl\'e and Masser)]
Assume that $A,B,C \in \mathbb{Z}$ are pairwise prime and satisfy the equation
$A+B=C$. Given $\epsilon>0$, there is a constant $M_\epsilon \in \mathbb{N}$
such that
\[
  \max\{|A|,|B|,|C|\}
    \leq \left(\prod_{p\mid ABC}p\right)^{1+\epsilon},
\]
where $p$ in the above equation is positive and prime.
\end{conjecture}

\begin{center}13\end{center}

\end{document}
